\documentclass[10pt,a4paper]{amsart}
\usepackage[margin=19mm]{geometry}
\usepackage{amsmath,amssymb,mathtools}
\usepackage[scr=rsfs,cal=euler]{mathalfa}
\usepackage{xfrac}
\usepackage[unicode,hidelinks,bookmarks=false]{hyperref}
\newcommand{\ie}{i.e.\ }
\newcommand{\cf}{cf.\ }
\newcommand{\resp}{respectively\ }
\newcommand{\Subsection}{\subsection}
\providecommand{\nobreakdash}{\nobreak\hbox{-}}
\setlength{\emergencystretch}{2em}
\multlinegap0pt
\usepackage[shortlabels]{enumitem}
\usepackage{amssymb}
\usepackage{xcolor}
\usepackage{tikz}
\usepackage{float}
\newtheorem{definition}{Definition}
\newtheorem{lemma}{Lemma}
\newtheorem{proposition}{Proposition}
\title{Homology of graph burnings}
\author{Yuri Muranov, Anna Muranova}
\date{Source: arXiv:2407.03832v1 (CC BY 4.0)}
\begin{document}
\maketitle

\noindent\emph{Source and licence.} Homology of graph burnings. Version arXiv:2407.03832v1 (CC BY 4.0), distributed by arXiv under the Creative Commons Attribution 4.0 International licence, http://creativecommons.org/licenses/by/4.0/, as recorded in the arXiv licence metadata for this exact version and shown on the abstract page https://arxiv.org/abs/2407.03832v1. Full licence notice: \texttt{licenses/arxiv-2407.03832-CC-BY-4.0.txt}.\par\smallskip

\noindent\textit{Editorial scope.} Counted unit: the article's proposition p3.11 (source0.tex lines 393--401). The article's complete original proof of it is delivered with it (source0.tex lines 402--462, one \texttt{proof} environment of 61 lines). No mathematics is written, repaired or re-derived: the statement and the proof are the article's own lines, in the article's own notation, and no retained line is substituted or re-flowed.  Retained uncounted as the source-local context the proof needs: lines 211--219; lines 223--227; lines 236--238; lines 241--243.  Nothing was omitted from any retained span, so the package carries no omission record.  No retained line contains a citation, so the wrapper carries no bibliography and no bibliography entry.  No retained line contains a figure environment, an image-inclusion command or drawing code, so no image is delivered and the package declares no asset dependency.  Editorial formatting only, in addition: an AMS \texttt{amsart} preamble that restates the article's own declarations and macros used by the retained lines, namely the environments \texttt{definition}, \texttt{lemma}, \texttt{proposition} on separate counters, together with the standard packages those lines need.  The retained environments print in this document as Definition 1, Lemma 1, Proposition 1 in order of appearance, whereas the article prints the same items with the numbering of its own full document.  The complete native source of the article as distributed in the arXiv e-print for this exact version is bundled unchanged as \texttt{sources/161821/source0.tex}.
\begin{equation}\label{3.3}
\begin{matrix}
\mathcal N_1 &\sqsubset & \mathcal N_2& \sqsubset &\mathcal N_3&\sqsubset \dots \sqsubset  &\mathcal N_k&\sqsubset&\mathcal N_{k+1}&\sqsubset  G\\
 && \sqcup&  &\sqcup &\dots & \sqcup&&||&\  \ \ || \\
                  &  &U_2&\sqsubset& U_3&\sqsubset\dots \sqsubset  &U_k&\sqsubset&U_{k+1}&\sqsubset  G\\
&& \sqcup&  &\sqcup &\dots & \sqcup&&\sqcup&\ \ \  \sqcup \\
                  &  &\mathcal N_1&\sqsubset& \mathcal N_2&\sqsubset\dots \sqsubset  &\mathcal N_{k-1}&\sqsubset&\mathcal N_{k}&\sqsubset  \mathcal N_{k+1}\\
\end{matrix}
\end{equation}

\begin{definition}\label{d3.1} \rm Let $G=(V,E)$ be a graph and 
$S_G=(v_1, \dots, v_k)$ be a non-empty ordered  set of vertices. 
 The sequence $S_G$ is called a \emph{burning sequence}  if  in  diagram (\ref{3.3}) $v_j \notin U_j$ for $2\leq j\leq k$ and  $U_{k+1}=G$. 
 The elements of $S_G$ are called \emph{burning sources}. The vertex $v_i$ is called the \emph{source at time $i$}. 
\end{definition}

\begin{equation}\label{3.5} 
\lambda_G\colon V_G\to \mathbb N\;\mbox{ by }\; \lambda_G(v)=\min\{i \,| \, v\in \mathcal N_i\}.
\end{equation}

\begin{lemma}\label{l3.2} Let $S_G=(v_1, \dots, v_n)$ be a burning sequence of a graph $G=(V,E)$.  Let $T=\max\{\lambda_G(v)\;|\;  \, v\in V_G\}$ where $\lambda_G$ is defined by \eqref{3.5}. Then the function $\lambda_G$ is a surjective map from  the set $V$ of vertices to the set  $\{1,2,\dots, T\}$,  
where  $T=k$ or $T=k+1$. 
\end{lemma} 

\begin{proposition}\label{p3.11} Fix a natural number $T\ge 1$ and consider burnings 
$\mathbf B^T(P_n)=(\lambda, S_{P_n})$ of  path graphs $P_n$ with the end time $T$.
\begin{enumerate}[label=$(\roman*)$]

\item  Then the maximum value of $n$  equals  $T^2$, i. e.  the longest path graph which can be burned  for the end time $T$ has $T^2$ vertices.  
\item 
If in additional $\lambda$ is a homomorphism, then  the maximum value of $n$  equals  $T^2-T+1$.
\end{enumerate}
\end{proposition}

\begin{proof}
(i) Firstly, we prove that  $n\le T^2$. By  Lemma \ref{l3.2}, we conclude 
that there are two possibilities for a burning  of a graph $P_n$ with  the end time $T$:

(1) a  burning of $P_n$ is given by a source sequence $S_{P_n}=(v_1,\dots, v_T)$ and 
\begin{equation}\label{3.6}
 \mathcal N_T= N_{T-1}(v_1)\Cup \dots \Cup N_{1}(v_{T-1})\Cup \mathcal N_0(v_T)=P_n,
\end{equation}
or, otherwise,  

(2)  a burning of $P_n$ is given by a source sequence $S_{P_n}=(v_1,\dots v_{T-1})$ and 
\begin{equation}\label{3.7}
 \ \mathcal N_{T}=N_{T-1}(v_1)\Cup  \dots \Cup N_{1}(v_{T-1})=P_n. \ \ \ \ \ \ \ \ 
\end{equation}
In the case (1), to obtain a maximum value $n$ we can take all the neighborhoods  in (\ref{3.6}) in such a way that pairwise distinct  neighborhoods in (\ref{3.6}) have empty intersection.  For a vertex  $v\in P_n$,  the maximal neighborhood 
$N_j(v)$  contains $2j+1$ vertices.  Hence, in the case (1),
\begin{equation}\label{3.8}
n\le \sum_{j=0}^{T-1}(2j+1)=T^2.
\end{equation}
In the case (2), using the same line of arguments   we obtain  that $n\le~T^2-1$.  Thus in both cases  $n\le T^2$. Finally, the burning sequence with $v_1=T$ and 
$$
v_j=\sum_{i=T-(j-1)}^{T-1}(2 i+1)+T-j+1 \ \ \text{for} \ \ 2\le j \le T,
$$
 defines the burning of the path graph  $P_{T^2}$.

(ii)   Now let $\lambda$ be a homomorphism. We consider this case similarly to the case (i).  At first we prove that the  equality (\ref{3.6}) leads to a contradiction. Note, that $n>1$, otherwise the statement is trivial. Since path graph $P_n$ is connected there is a vertex $v\in V_{P_n}$ such that $\{v, v_T\} \in E_{P_n}$. Moreover, $\lambda(v)=T-1$ since $\lambda$ is a homomorphism and $T$ is the maximal value of the map $\lambda$ on the set of vertices.  By Definition \ref{d3.1} we obtain
\begin{equation*}
v\in \mathcal N_{T-1}=N_{T-2}(v_1)\Cup N_{T-3}(v_2)\Cup \dots \Cup N_0(v_{T-1}). 
\end{equation*}
Thus, $v$ is a vertex of at least one of the graphs 
$N_{T-2}(v_1), \dots , N_0(v_{T-1})$.
  Since $d(v,v_T)=1$ we obtain that the vertex $v_T$ 
is the vertex of at least one of  the following  graphs 
$
N_{T-1}(v_1), \dots , N_1(v_{T-1})
$
and, hence, 
\begin{equation*}
v_T\in N_{T-1}(v_1)\Cup  N_{T-2}(v_2) \Cup \dots \Cup N_1(v_{T-1})= U_{T}. 
\end{equation*}
But, by Definition \ref{d3.1} and our assumption, $v_T\notin U_T$. Hence, we have obtained a contradiction in the case of the relation  (\ref{3.6}). 


Thus,  we can consider only the relation  (\ref{3.7}). 
The vertices  of the set $S_{P_n}=\{v_1, \dots , v_{T-1}\}$ are pairwise distinct  elements of the set of vertices 
$V_{P_{n}}=\{1, \dots , n\}$ and hence they  have a natural  order "$<$". Thus we obtain a sequence of  vertices $w_i \in V_{P_n} \ (1\leq i\leq T-1)$ and a permutation  $\pi$  of elements $(1, \dots, T-1)$ such that 
\begin{equation}\label{3.9}
1\leq w_1< w_2< \dots <w_{T-1}\leq n     \ \ \ \text{and} \ \  w_i=v_{\pi(j)}. 
\end{equation} 
For any two consequent vertices $w_i< w_{i+1}$ in (\ref{3.9}),  we 
have two neighborhoods from (\ref{3.7})
$M_{i}=N_{T-\pi^{-1}(i)}(w_i)$  and $M_{i+1}=N_{T-\pi^{-1}(i+1)}(w_{i+1})$.   To obtain   a maximum value of $n$ from above,  these neighborhoods must have minimal intersections.  The case of $M_{i}\cap M_{i+1}=\emptyset$ is impossible since in this case the burning map $\lambda$  is not a homomorphism.   Let  $M_i\cap M_{i+1}$ consists from  exactly one vertex for $1\leq i\leq T-2$. Then  the neighborhoods in  (\ref{3.7}) have $T-2$ one-vertex intersections since  the number of elements of this cover is $T-1$.  Hence, by (2) of proof  (i), we obtain 
\begin{equation}\label{3.10}
n\le T^2-1-(T-2)=T^2-T+1. 
\end{equation}
In this case, the burning sequence with $v_1=T$ and 
$$
v_j=\sum_{i=T-(j-1)}^{T-1}(2 i)+T-j  \ \ \text{for} \ \  j\ge 2
$$
 defines the burning of the path graph  $P_{T^2-T+1}$.
\end{proof}

\end{document}
